\section{Conclusion}
\label{sec:conclusion}

We identified a structural vulnerability in memory-based multi-modal fusion with foundation models: when modalities are processed sequentially through SAM2's memory mechanism, degraded sensors (\eg, RGB at night) pollute the shared memory bank and corrupt cross-modal attention for subsequent, otherwise reliable modalities.
To address this, we proposed Quality-Aware Memory Gating, a teacher--student framework that spatially modulates memory features before they are stored.
The teacher leverages SAM2's own decoder to assess per-modality spatial quality during training, while the lightweight student network (${\sim}$12.5K parameters) learns to predict quality maps at inference with negligible overhead.
Combined with Soft-MoE LoRA for modality-adaptive encoding and night-simulation augmentations for domain generalization, our approach addresses the key challenge of multi-modal maritime segmentation under varying sensor conditions.
Experiments on the MULTIAQUA benchmark validate our design choices.
